\documentclass[11pt]{article}

\usepackage[a4paper,margin=1in]{geometry}
\usepackage[T1]{fontenc}
\usepackage[utf8]{inputenc}
\usepackage{amsmath,amssymb,amsfonts,mathtools}
\usepackage{bm}
\usepackage{booktabs}
\usepackage{tabularx}
\usepackage{array}
\usepackage{algorithm}
\usepackage[noend]{algpseudocode}
\usepackage{hyperref}

\hypersetup{
    colorlinks=true,
    linkcolor=blue,
    citecolor=blue,
    urlcolor=blue
}

\newcommand{\Kset}{\mathcal{K}}
\newcommand{\Aset}{\mathcal{A}}
\newcommand{\Qset}{\mathcal{Q}}
\newcommand{\C}{\mathbb{C}}
\newcommand{\I}{\mathbf{I}}
\newcommand{\Hc}{\mathbf{H}}
\newcommand{\F}{\mathbf{F}}
\newcommand{\A}{\mathbf{A}}
\newcommand{\Sig}{\mathbf{\Sigma}}
\newcommand{\Smat}{\mathbf{S}}
\newcommand{\Wmat}{\mathbf{W}}
\newcommand{\tr}{\operatorname{tr}}
\newcommand{\norm}[1]{\left\lVert #1 \right\rVert}
\newcommand{\floor}[1]{\left\lfloor #1 \right\rfloor}
\newcommand{\UL}{\mathrm{UL}}
\newcommand{\DL}{\mathrm{DL}}
\newcommand{\E}{\mathbb{E}}
\emergencystretch=1.5em
\allowdisplaybreaks

\title{Finite-Blocklength Latency Minimization for Uplink and Downlink Communications:\\Theoretical Formulation and Solution Methods}
\author{}
\date{\today}

\begin{document}

\maketitle

\begin{abstract}
This report presents a theory-first account of the current uplink and downlink
work on finite-blocklength latency minimization. The exposition is organized as
a research document rather than an implementation note. It introduces a common
notation layer, derives the finite-blocklength rate expressions used in both
links, formulates the uplink and downlink latency-minimization problems under
payload-completion and fixed continuous communication scenarios, and then
develops two solution frameworks: the \emph{Training Only Method} and the
\emph{Training + Testing Method}. The emphasis is entirely on system models,
equations, constraints, and algorithms. Numerical configurations,
implementation-specific discussion, and simulation results are intentionally
excluded.
\end{abstract}

\section{Introduction}

Short-packet communication cannot be analyzed accurately through asymptotic
capacity alone. When the codeword length is finite, the achievable rate depends
not only on channel quality and transmit power, but also on blocklength and the
required decoding reliability. For latency-sensitive multiuser communication,
this creates a coupled design problem: the number of transmitted bits, the
chosen blocklength, and the precoder must be selected jointly.

The present work studies this problem in both uplink and downlink settings. In
the uplink, each user is handled through a user-wise sequence of transmission
blocks, and the main difficulty is that the number of required blocks is not
known a priori. In the downlink, all active users share the same transmission
epoch, so the main difficulty is inter-user interference and the resulting
coupling of feasible rates.

The report has four goals:
\begin{enumerate}
    \item establish a common mathematical notation for uplink and downlink,
    \item derive the corresponding finite-blocklength formulations,
    \item present the \emph{Training Only Method} and the \emph{Training + Testing Method} in algorithmic form, and
    \item define the communication scenarios used to interpret these formulations.
\end{enumerate}

The document is intentionally theoretical. It does not discuss software
structure, parameter files, implementation branches, or numerical result
artifacts. Instead, it focuses on the mathematical content of the current
methods.

\section{System Model and Common Notation}

\subsection{Core Variables}

Table~\ref{tab:common-notation} lists the symbols that are shared throughout
the report.

\begin{table}[ht]
\centering
\small
\caption{Common notation}
\label{tab:common-notation}
\renewcommand{\arraystretch}{1.15}
\begin{tabularx}{\textwidth}{>{\raggedright\arraybackslash}p{0.27\textwidth} X}
\toprule
Symbol & Meaning \\
\midrule
$\Kset = \{1,2,\dots,K\}$ & Set of users. \\
$k \in \Kset$ & User index. \\
$\ell$ & Transmission-block index. In uplink it is user-local; in downlink it is a common scheduling-block index. \\
$B_k^{\mathrm{tot}}$ & Total payload of user $k$ in bits. \\
$d_{k,\ell}$ & Requested bits for user $k$ in block $\ell$. This is the blockwise service demand presented to the optimizer. \\
$b_{k,\ell}$ & Bits actually served to user $k$ in block $\ell$. \\
$r_{k,\ell}$ & Remaining unsent bits of user $k$ when block $\ell$ starts in payload-completion mode. \\
$\bar b_{k,\ell}$ & Prescribed blockwise target of user $k$ in fixed continuous communication mode. \\
$n_{k,\ell}$ & Blocklength assigned to user $k$ in block $\ell$, measured in channel uses. \\
$T_k$ & Maximum admissible blocklength for user $k$ within one coherence-limited transmission opportunity. \\
$f_{s,k}$ & Symbol rate of user $k$ in symbols per second. \\
$\tau_k$ & End-to-end latency of user $k$. \\
$\varepsilon_k$ & Target packet error probability of user $k$. \\
$P_k$ & Uplink power budget of user $k$. \\
$P_B$ & Total transmit-power budget of the downlink block. \\
$d_k$ & Number of transmitted spatial streams of user $k$. \\
$\Aset_\ell$ & Active-user set in downlink block $\ell$. \\
$a_{k,\ell}$ & Anchor bits used in rollout-based training. These are already committed bits that must remain feasible at the queried blocklength. \\
\bottomrule
\end{tabularx}
\end{table}

The latency of user $k$ is defined by
\begin{equation}
\tau_k
=
\frac{1}{f_{s,k}}
\sum_{\ell} n_{k,\ell},
\label{eq:latency-user}
\end{equation}
and the network objective is to minimize the aggregate latency
\begin{equation}
\tau_{\mathrm{sum}}
=
\sum_{k\in\Kset}\tau_k.
\label{eq:latency-sum}
\end{equation}

\subsection{Finite-Blocklength Template}

For both links, the achievable rate is described through the normal
approximation \cite{polyanskiy2010}. Let mode $m \in \{\UL,\DL\}$ denote the
link direction. For user $k$ in block $\ell$, define the desired-signal
covariance $\Smat_{k,\ell}^{(m)}$ and the disturbance covariance
$\Wmat_{k,\ell}^{(m)}$. The whitened effective channel matrix is
\begin{equation}
\A_{k,\ell}^{(m)}
=
\bigl(\Wmat_{k,\ell}^{(m)}\bigr)^{-1/2}
\Smat_{k,\ell}^{(m)}
\bigl(\Wmat_{k,\ell}^{(m)}\bigr)^{-1/2}.
\label{eq:common-A}
\end{equation}

The corresponding capacity term is
\begin{equation}
C_{k,\ell}^{(m)}
=
\log_2\det\!\bigl(\I+\A_{k,\ell}^{(m)}\bigr),
\label{eq:common-C}
\end{equation}
and the channel-dispersion term is
\begin{equation}
V_{k,\ell}^{(m)}
=
(\log_2 e)^2
\tr\!\left[
\A_{k,\ell}^{(m)}
\bigl(\A_{k,\ell}^{(m)}+2\I\bigr)
\bigl(\A_{k,\ell}^{(m)}+\I\bigr)^{-2}
\right].
\label{eq:common-V}
\end{equation}

Therefore, for blocklength $n_{k,\ell}$ and reliability target
$\varepsilon_k$, the finite-blocklength rate is
\begin{equation}
R_{k,\ell}^{(m)}
=
C_{k,\ell}^{(m)}
-
\sqrt{\frac{V_{k,\ell}^{(m)}}{n_{k,\ell}}}\,
Q^{-1}(\varepsilon_k).
\label{eq:common-rate}
\end{equation}

If $b_{k,\ell}$ information bits are sent over $n_{k,\ell}$ channel uses,
feasibility requires
\begin{equation}
\frac{b_{k,\ell}}{n_{k,\ell}}
\le
R_{k,\ell}^{(m)}.
\label{eq:common-feasibility-rate}
\end{equation}
Multiplying both sides by $n_{k,\ell}$ yields
\begin{equation}
b_{k,\ell}
\le
n_{k,\ell}R_{k,\ell}^{(m)}.
\label{eq:common-feasibility-bits}
\end{equation}
Since the service variable is integer-valued, the largest feasible blockwise
service is
\begin{equation}
b_{k,\ell}^{\max,(m)}
=
\floor{n_{k,\ell}R_{k,\ell}^{(m)}}.
\label{eq:common-bmax}
\end{equation}

Equation~\eqref{eq:common-bmax} is central throughout the report: it converts a
continuous finite-blocklength rate into an integer blockwise service limit.

\section{Uplink Formulation}

\subsection{Effective Uplink Signal Model}

In the uplink, user $k$ in block $\ell$ is represented through the effective
MIMO relation
\begin{equation}
\mathbf y_{k,\ell}^{(\UL)}
=
\Hc_{k,\ell}^{(\UL)}
\F_{k,\ell}^{(\UL)}
\mathbf s_{k,\ell}
+
\mathbf z_{k,\ell}^{(\UL)},
\label{eq:ul-signal}
\end{equation}
where
\begin{equation}
\Hc_{k,\ell}^{(\UL)} \in \C^{N_{r,k}\times N_{t,k}},
\qquad
\F_{k,\ell}^{(\UL)} \in \C^{N_{t,k}\times d_k},
\label{eq:ul-dimensions}
\end{equation}
$\mathbf s_{k,\ell}$ is the transmitted symbol vector with identity covariance,
and $\mathbf z_{k,\ell}^{(\UL)}$ is the effective disturbance term with
covariance $\Wmat_{k,\ell}^{(\UL)}$.

The desired-signal covariance is
\begin{equation}
\Smat_{k,\ell}^{(\UL)}
=
\Hc_{k,\ell}^{(\UL)}
\F_{k,\ell}^{(\UL)}
\bigl(\F_{k,\ell}^{(\UL)}\bigr)^H
\bigl(\Hc_{k,\ell}^{(\UL)}\bigr)^H,
\label{eq:ul-signal-cov}
\end{equation}
and the whitened matrix becomes
\begin{equation}
\A_{k,\ell}^{(\UL)}
=
\bigl(\Wmat_{k,\ell}^{(\UL)}\bigr)^{-1/2}
\Smat_{k,\ell}^{(\UL)}
\bigl(\Wmat_{k,\ell}^{(\UL)}\bigr)^{-1/2}.
\label{eq:ul-A}
\end{equation}
Substituting \eqref{eq:ul-A} into
\eqref{eq:common-C}--\eqref{eq:common-rate} gives
$C_{k,\ell}^{(\UL)}$, $V_{k,\ell}^{(\UL)}$, and
$R_{k,\ell}^{(\UL)}$.

\subsection{Single-Block Uplink Design Problem}

Suppose block $\ell$ of user $k$ is asked to serve $d_{k,\ell}$ bits. The
minimum-latency one-block design problem is
\begin{subequations}
\begin{align}
\min_{n_{k,\ell},\,\F_{k,\ell}^{(\UL)}}
\quad
&
\frac{n_{k,\ell}}{f_{s,k}}
\label{eq:ul-single-obj}
\\
\text{s.t.}\quad
&
\frac{d_{k,\ell}}{n_{k,\ell}}
\le
R_{k,\ell}^{(\UL)},
\label{eq:ul-single-rate}
\\
&
1 \le n_{k,\ell} \le T_k,
\label{eq:ul-single-n}
\\
&
\norm{\F_{k,\ell}^{(\UL)}}_F^2 \le P_k.
\label{eq:ul-single-power}
\end{align}
\end{subequations}

Problem \eqref{eq:ul-single-obj}--\eqref{eq:ul-single-power} can also be read
from the dual viewpoint. For fixed blocklength, the objective is to enlarge the
finite-blocklength rate until the requested service becomes feasible; for fixed
precoder, the objective is to find the smallest admissible blocklength.

\subsection{Uplink Payload-Completion Formulation}

In payload completion mode, the number of blocks required by user $k$ is not
known in advance. Let $L_k$ denote the number of blocks eventually used by user
$k$. The uplink payload problem is
\begin{subequations}
\begin{align}
\min_{\{L_k,b_{k,\ell},n_{k,\ell},\F_{k,\ell}^{(\UL)}\}}
\quad
&
\sum_{k\in\Kset}\frac{1}{f_{s,k}}\sum_{\ell=1}^{L_k} n_{k,\ell}
\label{eq:ul-payload-obj}
\\
\text{s.t.}\quad
&
\sum_{\ell=1}^{L_k} b_{k,\ell} = B_k^{\mathrm{tot}},
\qquad \forall k,
\label{eq:ul-payload-conservation}
\\
&
b_{k,\ell}
\le
\floor{n_{k,\ell}R_{k,\ell}^{(\UL)}},
\qquad \forall k,\ell,
\label{eq:ul-payload-fbl}
\\
&
1 \le n_{k,\ell} \le T_k,
\qquad \forall k,\ell,
\label{eq:ul-payload-n}
\\
&
\norm{\F_{k,\ell}^{(\UL)}}_F^2 \le P_k,
\qquad \forall k,\ell.
\label{eq:ul-payload-power}
\end{align}
\end{subequations}

The remaining-bits recursion is
\begin{equation}
r_{k,1}=B_k^{\mathrm{tot}},
\qquad
r_{k,\ell+1}=r_{k,\ell}-b_{k,\ell},
\label{eq:ul-remaining-recursion}
\end{equation}
and the scheduler terminates once $r_{k,\ell}=0$.

\subsection{Uplink Fixed Continuous Communication Formulation}

In fixed continuous communication mode, the horizon length is prescribed in
advance, and each block carries a target $\bar b_{k,\ell}$. The problem becomes
\begin{subequations}
\begin{align}
\min_{\{b_{k,\ell},n_{k,\ell},\F_{k,\ell}^{(\UL)}\}}
\quad
&
\sum_{k\in\Kset}\frac{1}{f_{s,k}}\sum_{\ell=1}^{L} n_{k,\ell}
\label{eq:ul-fixed-obj}
\\
\text{s.t.}\quad
&
0 \le b_{k,\ell} \le \bar b_{k,\ell},
\qquad \forall k,\ell,
\label{eq:ul-fixed-target}
\\
&
b_{k,\ell}
\le
\floor{n_{k,\ell}R_{k,\ell}^{(\UL)}},
\qquad \forall k,\ell,
\label{eq:ul-fixed-fbl}
\\
&
1 \le n_{k,\ell} \le T_k,
\qquad
\norm{\F_{k,\ell}^{(\UL)}}_F^2 \le P_k.
\label{eq:ul-fixed-other}
\end{align}
\end{subequations}

The key difference from payload completion is that the horizon length $L$ is
fixed and blockwise shortfall
$\bar b_{k,\ell}-b_{k,\ell}$ is recorded explicitly rather than being carried
forward as an unscheduled residual payload.

\section{Downlink Formulation}

\subsection{Interference-Coupled Downlink Signal Model}

In the downlink, all active users in block $\ell$ are transmitted
simultaneously. The received signal of user $k$ is
\begin{equation}
\mathbf y_{k,\ell}^{(\DL)}
=
\Hc_{k,\ell}^{(\DL)}\F_{k,\ell}^{(\DL)}\mathbf s_{k,\ell}
+
\sum_{j\in\Aset_\ell\setminus\{k\}}
\Hc_{k,\ell}^{(\DL)}\F_{j,\ell}^{(\DL)}\mathbf s_{j,\ell}
+
\mathbf z_{k,\ell}^{(\DL)},
\label{eq:dl-signal}
\end{equation}
where
\begin{equation}
\Hc_{k,\ell}^{(\DL)} \in \C^{N_{r,k}\times N_b},
\qquad
\F_{k,\ell}^{(\DL)} \in \C^{N_b\times d_k},
\label{eq:dl-dimensions}
\end{equation}
and $\mathbf z_{k,\ell}^{(\DL)}$ is receiver noise.

The desired-signal covariance of user $k$ is
\begin{equation}
\Smat_{k,\ell}^{(\DL)}
=
\Hc_{k,\ell}^{(\DL)}
\F_{k,\ell}^{(\DL)}
\bigl(\F_{k,\ell}^{(\DL)}\bigr)^H
\bigl(\Hc_{k,\ell}^{(\DL)}\bigr)^H.
\label{eq:dl-signal-cov}
\end{equation}
Its disturbance covariance is
\begin{equation}
\Wmat_{k,\ell}^{(\DL)}
=
\sigma_k^2\I
+
\sum_{j\in\Aset_\ell\setminus\{k\}}
\Hc_{k,\ell}^{(\DL)}
\F_{j,\ell}^{(\DL)}
\bigl(\F_{j,\ell}^{(\DL)}\bigr)^H
\bigl(\Hc_{k,\ell}^{(\DL)}\bigr)^H.
\label{eq:dl-disturbance-cov}
\end{equation}

Hence the downlink whitened matrix is
\begin{equation}
\A_{k,\ell}^{(\DL)}
=
\bigl(\Wmat_{k,\ell}^{(\DL)}\bigr)^{-1/2}
\Smat_{k,\ell}^{(\DL)}
\bigl(\Wmat_{k,\ell}^{(\DL)}\bigr)^{-1/2},
\label{eq:dl-A}
\end{equation}
which leads directly to
$C_{k,\ell}^{(\DL)}$, $V_{k,\ell}^{(\DL)}$, and
$R_{k,\ell}^{(\DL)}$ through
\eqref{eq:common-C}--\eqref{eq:common-rate}.

\subsection{Per-Block Downlink Design Problem}

Let $\Aset_\ell$ be the active-user set in block $\ell$, and let $d_{k,\ell}$
be the requested bits of active user $k$. The one-block downlink problem is
\begin{subequations}
\begin{align}
\min_{\{n_{k,\ell},\F_{k,\ell}^{(\DL)}\}_{k\in\Aset_\ell}}
\quad
&
\sum_{k\in\Aset_\ell}\frac{n_{k,\ell}}{f_{s,k}}
\label{eq:dl-single-obj}
\\
\text{s.t.}\quad
&
\frac{d_{k,\ell}}{n_{k,\ell}}
\le
R_{k,\ell}^{(\DL)},
\qquad \forall k\in\Aset_\ell,
\label{eq:dl-single-rate}
\\
&
1 \le n_{k,\ell} \le T_k,
\qquad \forall k\in\Aset_\ell,
\label{eq:dl-single-n}
\\
&
\sum_{k\in\Aset_\ell}\norm{\F_{k,\ell}^{(\DL)}}_F^2 \le P_B.
\label{eq:dl-single-power}
\end{align}
\end{subequations}

The complication relative to the uplink is immediate:
$R_{k,\ell}^{(\DL)}$ depends on the entire active-user beam set through
\eqref{eq:dl-disturbance-cov}. Thus, even if only one user's precoder is
changed, the feasible rates of all other active users are altered as well.

\subsection{Downlink Payload-Completion Formulation}

In payload completion mode, the active-user set is induced by the remaining
payload:
\begin{equation}
\Aset_\ell
=
\left\{
k\in\Kset : r_{k,\ell}>0
\right\}.
\label{eq:dl-active-payload}
\end{equation}

The network-level problem is
\begin{subequations}
\begin{align}
\min_{\{b_{k,\ell},n_{k,\ell},\F_{k,\ell}^{(\DL)}\}}
\quad
&
\sum_{k\in\Kset}\frac{1}{f_{s,k}}\sum_{\ell} n_{k,\ell}
\label{eq:dl-payload-obj}
\\
\text{s.t.}\quad
&
\sum_{\ell} b_{k,\ell} = B_k^{\mathrm{tot}},
\qquad \forall k,
\label{eq:dl-payload-conservation}
\\
&
b_{k,\ell}
\le
\floor{n_{k,\ell}R_{k,\ell}^{(\DL)}},
\qquad \forall k,\ell,
\label{eq:dl-payload-fbl}
\\
&
1 \le n_{k,\ell} \le T_k,
\qquad \forall k,\ell,
\label{eq:dl-payload-n}
\\
&
\sum_{k\in\Aset_\ell}\norm{\F_{k,\ell}^{(\DL)}}_F^2 \le P_B,
\qquad \forall \ell.
\label{eq:dl-payload-power}
\end{align}
\end{subequations}

The same remaining-bits recursion as in
\eqref{eq:ul-remaining-recursion} applies, but now the block service of one
user depends on the simultaneous service decisions of the others.

\subsection{Downlink Fixed Continuous Communication Formulation}

In fixed continuous communication mode, the block horizon is prescribed and the
blockwise targets are fixed:
\begin{subequations}
\begin{align}
\min_{\{b_{k,\ell},n_{k,\ell},\F_{k,\ell}^{(\DL)}\}}
\quad
&
\sum_{k\in\Kset}\frac{1}{f_{s,k}}\sum_{\ell=1}^{L} n_{k,\ell}
\label{eq:dl-fixed-obj}
\\
\text{s.t.}\quad
&
0 \le b_{k,\ell} \le \bar b_{k,\ell},
\qquad \forall k,\ell,
\label{eq:dl-fixed-target}
\\
&
b_{k,\ell}
\le
\floor{n_{k,\ell}R_{k,\ell}^{(\DL)}},
\qquad \forall k,\ell,
\label{eq:dl-fixed-fbl}
\\
&
1 \le n_{k,\ell} \le T_k,
\qquad \forall k,\ell,
\label{eq:dl-fixed-n}
\\
&
\sum_{k\in\Aset_\ell}\norm{\F_{k,\ell}^{(\DL)}}_F^2 \le P_B,
\qquad \forall \ell.
\label{eq:dl-fixed-power}
\end{align}
\end{subequations}

The difference from payload completion is again conceptual: the demand sequence
is externally prescribed, so any shortfall
$\bar b_{k,\ell}-b_{k,\ell}$ is treated as unmet service in that block rather
than as a residual workload to be automatically rolled forward.

\section{Training Only Method}

\subsection{Common Primal-Dual Structure}

The \emph{Training Only Method} performs optimization directly on the current
channel realization. There is no offline learning stage. Instead, each decision
block is solved online through a constrained primal-dual update.

For uplink block $(k,\ell)$ with requested bits $d_{k,\ell}$, define the rate
constraint residual
\begin{equation}
g_{k,\ell}^{(\UL,r)}
=
\frac{d_{k,\ell}}{n_{k,\ell}} - R_{k,\ell}^{(\UL)},
\label{eq:ul-rate-residual}
\end{equation}
and the power residual
\begin{equation}
g_{k,\ell}^{(\UL,p)}
=
\norm{\F_{k,\ell}^{(\UL)}}_F^2 - P_k.
\label{eq:ul-power-residual}
\end{equation}

For the downlink block, define
\begin{equation}
g_{k,\ell}^{(\DL,r)}
=
\frac{d_{k,\ell}}{n_{k,\ell}} - R_{k,\ell}^{(\DL)},
\qquad k\in\Aset_\ell,
\label{eq:dl-rate-residual}
\end{equation}
and the common block-power residual
\begin{equation}
g_{\ell}^{(\DL,p)}
=
\sum_{k\in\Aset_\ell}\norm{\F_{k,\ell}^{(\DL)}}_F^2 - P_B.
\label{eq:dl-power-residual}
\end{equation}

The blockwise Lagrangian is then formed by combining a rate-maximizing primal
objective with nonnegative dual penalties. In uplink,
\begin{equation}
\mathcal{L}_{k,\ell}^{(\UL)}
=
-R_{k,\ell}^{(\UL)}
+
\lambda_{k,\ell}^{(\UL,r)}\bigl[g_{k,\ell}^{(\UL,r)}\bigr]_+
+
\lambda_{k,\ell}^{(\UL,p)}\bigl[g_{k,\ell}^{(\UL,p)}\bigr]_+,
\label{eq:ul-lagrangian}
\end{equation}
while in downlink,
\begin{equation}
\mathcal{L}_{\ell}^{(\DL)}
=
-\sum_{k\in\Aset_\ell} R_{k,\ell}^{(\DL)}
+
\sum_{k\in\Aset_\ell}
\lambda_{k,\ell}^{(\DL,r)}\bigl[g_{k,\ell}^{(\DL,r)}\bigr]_+
+
\lambda_{\ell}^{(\DL,p)}\bigl[g_{\ell}^{(\DL,p)}\bigr]_+.
\label{eq:dl-lagrangian}
\end{equation}

The primal variables are updated to reduce the Lagrangian, and the dual
variables are updated to penalize persistent infeasibility. To monitor the
solve, the method tracks three KKT-style residuals:
\begin{align}
r_p
&=
\text{largest positive primal violation},
\label{eq:kkt-primal}
\\
r_c
&=
\text{largest complementarity product between a multiplier and its positive violation},
\label{eq:kkt-complementarity}
\\
r_s
&=
\text{relative change of the optimized precoder state between successive updates}.
\label{eq:kkt-stationarity}
\end{align}
The online solve is terminated when these residuals satisfy prescribed
tolerances, or when the state becomes practically stationary while still
violating feasibility.

\subsection{Uplink Training Only Method}

The uplink training-only method is fundamentally sequential. A user first tries
to serve its current request with the largest available blocklength $T_k$. The
maximum full-block service is
\begin{equation}
b_{k,\ell}^{(T)}
=
\min\!\left\{
d_{k,\ell},
\floor{T_k\,R_{k,\ell}^{(\UL)}(T_k)}
\right\}.
\label{eq:ul-full-block-service}
\end{equation}

If $b_{k,\ell}^{(T)} < d_{k,\ell}$, the block is only partially served and the
method immediately opens a new block for the unsent remainder. If
$b_{k,\ell}^{(T)} = d_{k,\ell}$, then the same bit load is kept fixed while the
method attempts to reduce blocklength. At each candidate smaller blocklength,
the precoder is re-optimized and the candidate is accepted only if the same
committed bit count remains feasible.

\begin{algorithm}[ht]
\caption{Training-Only Uplink Scheduler for One User}
\label{alg:ul-train-only}
\begin{algorithmic}[1]
\Require Payload rule generating requests $\{d_{k,\ell}\}$ for user $k$
\Ensure Block sequence $\{(b_{k,\ell},n_{k,\ell},\F_{k,\ell}^{(\UL)})\}$
\State Initialize $\ell \gets 1$
\While{a new request $d_{k,\ell}$ exists}
    \State Solve the full-block problem \eqref{eq:ul-single-obj}--\eqref{eq:ul-single-power} at $n_{k,\ell}=T_k$ using the primal-dual updates of \eqref{eq:ul-lagrangian}
    \State Compute $b_{k,\ell}^{(T)}$ from \eqref{eq:ul-full-block-service}
    \If{$b_{k,\ell}^{(T)} = 0$}
        \State Mark the block as unserved and stop the current request
    \EndIf
    \If{$b_{k,\ell}^{(T)} < d_{k,\ell}$}
        \State Commit $b_{k,\ell}=b_{k,\ell}^{(T)}$ and set $n_{k,\ell}=T_k$
        \State Create a new request for the unsent remainder
    \Else
        \State Set $b_{k,\ell}=d_{k,\ell}$ and initialize $n_{k,\ell}=T_k$
        \While{a smaller candidate blocklength remains feasible after re-optimization}
            \State Accept the smaller blocklength and update the precoder
        \EndWhile
        \State Commit the smallest accepted $n_{k,\ell}$
    \EndIf
    \State $\ell \gets \ell + 1$
\EndWhile
\end{algorithmic}
\end{algorithm}

Algorithm~\ref{alg:ul-train-only} solves both uplink scenarios. In payload
completion mode, the next request is the remaining payload. In fixed continuous
communication mode, the next request is simply the next prescribed target.

\subsection{Downlink Training Only Method}

In the downlink, a whole active-user block must be solved jointly. Let
$d_{k,\ell}$ be the current request of user $k\in\Aset_\ell$. The training-only
method first performs block-level precoder convergence at the full coherence
limits $\{T_k\}_{k\in\Aset_\ell}$. This produces full-block supported service
\begin{equation}
b_{k,\ell}^{(T)}
=
\min\!\left\{
d_{k,\ell},
\floor{T_k\,R_{k,\ell}^{(\DL)}(T_k)}
\right\},
\qquad k\in\Aset_\ell.
\label{eq:dl-full-block-service}
\end{equation}

Users with $b_{k,\ell}^{(T)}=0$ are removed from transmission in that block. For
users that are fully served at $T_k$, the method tries to reduce blocklength
while keeping the already committed bits fixed. Since the block is coupled, a
candidate reduction may invalidate another user's committed service. The method
therefore re-optimizes the affected active-user subset whenever necessary and
accepts a candidate reduction only if all committed bits remain feasible.

\begin{algorithm}[ht]
\caption{Training-Only Downlink Block Solver}
\label{alg:dl-train-only}
\begin{algorithmic}[1]
\Require Active-user set $\Aset_\ell$ and requests $\{d_{k,\ell}\}_{k\in\Aset_\ell}$
\Ensure Block plan $\{(b_{k,\ell},n_{k,\ell},\F_{k,\ell}^{(\DL)})\}_{k\in\Aset_\ell}$
\State Initialize the active-user precoders
\Repeat
    \State Update the active-user precoders using the block Lagrangian \eqref{eq:dl-lagrangian}
    \State Update the multipliers $\{\lambda_{k,\ell}^{(\DL,r)}\}$ and $\lambda_{\ell}^{(\DL,p)}$
    \State Measure the residuals $r_p$, $r_c$, and $r_s$
\Until{KKT convergence or stationary infeasibility}
\State Restore the best feasible block state if one exists; otherwise restore the best minimum-violation state
\For{each active user $k\in\Aset_\ell$}
    \State Compute $b_{k,\ell}^{(T)}$ from \eqref{eq:dl-full-block-service}
\EndFor
\State Remove users with zero supported service from the transmit set of the block
\For{each user fully served at $T_k$}
    \While{a smaller candidate blocklength preserves all already committed bits}
        \State Accept the candidate directly if feasible
        \State Otherwise re-optimize the affected active-user subset and accept only if the committed bits of all active users remain feasible
    \EndWhile
\EndFor
\State Commit the final joint block plan
\end{algorithmic}
\end{algorithm}

The training-only downlink method is therefore a \emph{converge first, allocate
second, reduce third} procedure. This ordering is essential because the
finite-blocklength rate of each user depends on the full active-user beam set.

\section{Training + Testing Method}

\subsection{Method Principle}

The \emph{Training + Testing Method} separates optimization into two stages.
During training, a precoder mapping is learned from a collection of rollout
queries produced by the same finite-blocklength scheduling logic that will later
be used at test time. During testing, online gradient-based optimization is
replaced by model inference, followed by the same feasibility checks and the
same blocklength-reduction logic.

The central training object is a rollout query. In general form, a query
$q\in\Qset$ contains
\begin{equation}
q
=
\bigl(
\text{channel state},
\text{blocklength state},
\text{disturbance state},
\text{anchor bits},
\text{reliability target},
\text{phase label}
\bigr).
\label{eq:general-rollout-query}
\end{equation}

The phase label is one of the following:
\begin{enumerate}
    \item \textbf{full\_block}: the model is queried at the largest admissible blocklength;
    \item \textbf{tail\_feasible}: the model is queried at a reduced blocklength while the committed anchor bits remain feasible;
    \item \textbf{tail\_frontier}: the model is queried at the first rejected reduced blocklength beyond the last feasible point.
\end{enumerate}

The anchor bits play a crucial role. They transform the training problem from a
generic rate-maximization task into a \emph{boundary-learning} task: the model
is explicitly trained to distinguish feasible and nearly infeasible
finite-blocklength operating points.

\subsection{Uplink Training + Testing Method}

\subsubsection{Training stage}

For uplink user $k$, a rollout query has the abstract form
\begin{equation}
q_{k,\ell}^{(\UL)}
=
\bigl(
\Hc_{k,\ell}^{(\UL)},
\Wmat_{k,\ell}^{(\UL)},
n,
a_{k,\ell},
\varepsilon_k,
P_k,
\varphi_{k,\ell}
\bigr),
\label{eq:ul-query}
\end{equation}
where $\varphi_{k,\ell}$ is the phase label.

The rollout is generated block by block. For a request $d_{k,\ell}$:
\begin{enumerate}
    \item create a \texttt{full\_block} query at $n=T_k$ with anchor bits $a_{k,\ell}=0$;
    \item compute the full-block committed service
    \begin{equation}
    b_{k,\ell}^{\mathrm{com}}
    =
    \min\!\left\{
    d_{k,\ell},
    \floor{T_k\,R_{k,\ell}^{(\UL)}(T_k)}
    \right\};
    \label{eq:ul-committed-bits}
    \end{equation}
    \item if $b_{k,\ell}^{\mathrm{com}}<d_{k,\ell}$, the block is partially served and tail exploration stops for that block;
    \item if $b_{k,\ell}^{\mathrm{com}}=d_{k,\ell}$, reduce the blocklength and record every feasible reduced point as \texttt{tail\_feasible};
    \item record the first rejected reduced point as \texttt{tail\_frontier}.
\end{enumerate}

Let $\Theta_k$ denote the parameters of the user-$k$ model. A generic
rollout-weighted training objective is
\begin{equation}
\mathcal{J}_k^{(\UL)}(\Theta_k)
=
\E_{q\in\Qset_k^{(\UL)}}
\left[
\omega(q)
\left(
-R(q;\Theta_k)
+
\lambda_k^{(r)}\Bigl[\frac{a(q)}{n(q)}-R(q;\Theta_k)\Bigr]_+
+
\lambda_k^{(p)}\bigl[\Pi(q;\Theta_k)-P_k\bigr]_+
\right)
\right],
\label{eq:ul-training-objective}
\end{equation}
where $\omega(q)$ is the query weight, $a(q)$ is the anchor-bit count,
$n(q)$ is the queried blocklength, and $\Pi(q;\Theta_k)$ denotes the transmit
power induced by the predicted precoder.

\begin{algorithm}[ht]
\caption{Uplink Training + Testing Method}
\label{alg:ul-train-test}
\begin{algorithmic}[1]
\Require Training episodes and a testing scenario for user $k$
\Ensure Learned model and final test-time block plan
\State \textbf{Training rollout generation:}
\For{each training episode and each request block}
    \State Record one \texttt{full\_block} query at $n=T_k$
    \State Compute the committed bits $b_{k,\ell}^{\mathrm{com}}$ from \eqref{eq:ul-committed-bits}
    \If{$b_{k,\ell}^{\mathrm{com}} = d_{k,\ell}$}
        \State Decrease $n$ and record feasible reduced states as \texttt{tail\_feasible}
        \State Record the first rejected state as \texttt{tail\_frontier}
    \EndIf
\EndFor
\State Train the user model by minimizing \eqref{eq:ul-training-objective}
\State \textbf{Testing:}
\While{a new request $d_{k,\ell}$ exists}
    \State Infer the precoder at $n=T_k$
    \State Compute $b_{k,\ell}^{(T)}=\min\{d_{k,\ell},\floor{T_kR_{k,\ell}^{(\UL)}(T_k)}\}$
    \If{$b_{k,\ell}^{(T)}=d_{k,\ell}$}
        \State Reduce $n$ greedily while keeping $b_{k,\ell}^{(T)}$ fixed
    \EndIf
    \State Commit the smallest feasible block and move to the next request
\EndWhile
\end{algorithmic}
\end{algorithm}

\subsubsection{Testing stage}

At test time, the model replaces the online primal update, but the scheduling
logic remains unchanged. The full-block prediction is queried first, a feasible
bit count is committed, and blocklength reduction is performed only when the
entire current request is already supported.

\subsection{Downlink Training + Testing Method}

\subsubsection{Training stage}

The downlink training state is joint across users. A rollout state for block
$\ell$ can be written as
\begin{equation}
q_{\ell}^{(\DL)}
=
\bigl(
\{\Hc_{k,\ell}^{(\DL)}\}_{k\in\Kset},
\mathbf m_\ell,
\mathbf n_\ell,
\mathbf a_\ell,
\{\Wmat_{k,\ell}^{(\DL)}\}_{k\in\Kset},
\{\varepsilon_k\}_{k\in\Kset},
\varphi_\ell
\bigr),
\label{eq:dl-query}
\end{equation}
where $\mathbf m_\ell$ is the active or service mask, $\mathbf n_\ell$ is the
vector of candidate blocklengths, and $\mathbf a_\ell$ is the vector of anchor
bits.

The rollout is built causally:
\begin{enumerate}
    \item generate a \texttt{full\_block} query at $\mathbf n_\ell=\mathbf T$;
    \item compute the supported bits at the full block;
    \item form a \emph{service mask} that removes users with zero supported service;
    \item if the service mask differs from the active mask, record a second
    \texttt{full\_block} query for the masked state;
    \item define committed bits
    \begin{equation}
    b_{k,\ell}^{\mathrm{com}}
    =
    \min\!\left\{
    d_{k,\ell},
    \floor{T_k\,R_{k,\ell}^{(\DL)}(T_k)}
    \right\};
    \label{eq:dl-committed-bits}
    \end{equation}
    \item identify reducible users, namely those already fully served at the
    full block;
    \item generate \texttt{tail\_feasible} queries by decreasing the candidate
    blocklengths of reducible users while holding the committed bits fixed;
    \item record the first rejected reduced joint state as
    \texttt{tail\_frontier}.
\end{enumerate}

Let $\Theta$ denote the parameters of the joint downlink model collection. A
generic rollout-weighted objective is
\begin{equation}
\mathcal{J}^{(\DL)}(\Theta)
=
\E_{q\in\Qset^{(\DL)}}
\left[
\omega(q)
\left(
-\sum_{k\in\Aset(q)} R_k(q;\Theta)
+
\sum_{k\in\Aset(q)}
\lambda_k^{(r)}
\Bigl[\frac{a_k(q)}{n_k(q)}-R_k(q;\Theta)\Bigr]_+
+
\lambda^{(p)}
\Bigl[
\sum_{k\in\Aset(q)}\Pi_k(q;\Theta)-P_B
\Bigr]_+
\right)
\right].
\label{eq:dl-training-objective}
\end{equation}

\begin{algorithm}[ht]
\caption{Downlink Training + Testing Method}
\label{alg:dl-train-test}
\begin{algorithmic}[1]
\Require Training episodes and one testing scenario
\Ensure Learned model collection and final test-time schedule
\State \textbf{Training rollout generation:}
\For{each training block}
    \State Record a \texttt{full\_block} query at the full blocklength vector $\mathbf T$
    \State Remove zero-service users and, if necessary, record the masked \texttt{full\_block} state
    \State Compute committed bits from \eqref{eq:dl-committed-bits}
    \If{some users are fully served at the full block}
        \State Decrease their candidate blocklengths while the committed bits remain feasible
        \State Record feasible reduced states as \texttt{tail\_feasible}
        \State Record the first rejected reduced state as \texttt{tail\_frontier}
    \EndIf
\EndFor
\State Train the model collection by minimizing \eqref{eq:dl-training-objective}
\State \textbf{Testing:}
\For{each communication block}
    \State Determine the active-user set from the scenario state
    \State Infer the active-user precoders at the full blocklength vector
    \State Remove any zero-service users from the transmit set
    \State Commit bits at the full blocklength vector
    \State For each fully served user, reduce blocklength greedily while preserving all committed bits
    \State Re-evaluate the joint committed state and finalize the block plan
\EndFor
\end{algorithmic}
\end{algorithm}

\subsubsection{Testing stage}

The testing stage remains block-causal. The learned model is asked to provide
precoders for the currently active block state, after which the same service
masking, bit commitment, and blocklength-reduction logic is applied. Thus, the
learned component replaces online beam optimization, but not the finite-
blocklength feasibility structure of the scheduler itself.

\section{Communication Scenarios}

\subsection{Payload}

In payload-completion mode, each user begins with a finite payload
$B_k^{\mathrm{tot}}$. The scheduler stops only when this payload is fully
drained. The request presented to a block is the current remainder:
\begin{equation}
d_{k,\ell}=r_{k,\ell}.
\label{eq:payload-request}
\end{equation}
After the block is committed, the remainder is updated as
\begin{equation}
r_{k,\ell+1}=r_{k,\ell}-b_{k,\ell}.
\label{eq:payload-update}
\end{equation}

\subsection{Fixed Continuous Communication (Bits per Block)}

In fixed continuous communication mode, the horizon length $L$ is known in
advance and the blockwise targets are specified externally:
\begin{equation}
d_{k,\ell}=\bar b_{k,\ell},
\qquad \ell=1,\dots,L.
\label{eq:fixed-request}
\end{equation}
The unmet service in a block is
\begin{equation}
u_{k,\ell}=\bar b_{k,\ell}-b_{k,\ell},
\qquad
u_{k,\ell}\ge 0,
\label{eq:fixed-shortfall}
\end{equation}
which is recorded as blockwise shortfall rather than automatically rescheduled
as a new payload remainder.

\subsection{Uplink Payload}

For uplink payload completion, each user evolves independently according to its
own remainder recursion. User $k$ creates a new block only if its current block
cannot finish the entire remaining payload. Consequently, the block count
$L_k$ is an output of the optimization.

\subsection{Uplink Continuous Communication}

For uplink fixed continuous communication, the decision epochs are prescribed in
advance. Each epoch has a target $\bar b_{k,\ell}$, and the main question is
how much of that target can be supported within the coherence and power limits
of the user at that epoch.

\subsection{Downlink Payload}

For downlink payload completion, the active-user set changes with time:
\begin{equation}
\Aset_\ell = \{k\in\Kset : r_{k,\ell}>0\}.
\label{eq:downlink-payload-active}
\end{equation}
As users finish their payloads they leave the active set, which changes both
the interference structure and the feasible rate region of the remaining users.

\subsection{Downlink Continuous Communication}

For downlink fixed continuous communication, the block horizon is prescribed and
every block carries a target vector $\{\bar b_{k,\ell}\}_{k\in\Kset}$. The
active-user set is determined by the users that are scheduled in that block. If
the coupled downlink block cannot support a user's requested service, the user
may receive zero service or partial service in that block, and this outcome is
recorded as part of the schedule.

\section{Conclusion}

The report has recast the current work into a theory-centered structure. The
development proceeds from a common finite-blocklength model, to separate uplink
and downlink formulations, and then to two solution frameworks.

The uplink problem is characterized by user-wise sequential block creation under
finite-blocklength feasibility. The downlink problem is characterized by a
jointly coupled interference structure that requires block-level coordination
before blockwise service can be committed.

The \emph{Training Only Method} solves each decision state online through
constrained primal-dual optimization and feasibility-guided blocklength
reduction. The \emph{Training + Testing Method} replaces online optimization by
learned precoder inference, but retains the same finite-blocklength service
logic through rollout training, committed-bit anchors, and frontier states.

Taken together, these formulations provide a complete theoretical description of
the current uplink and downlink methodology without relying on implementation
details or numerical experiment summaries.

\begin{thebibliography}{1}

\bibitem{polyanskiy2010}
Y.~Polyanskiy, H.~V.~Poor, and S.~Verd\'u,
``Channel coding rate in the finite blocklength regime,''
\emph{IEEE Transactions on Information Theory},
vol.~56, no.~5, pp.~2307--2359, 2010.

\end{thebibliography}

\end{document}
